% 数学ノート集・共通テンプレート v2（積分論2026年度第1部の版面を参考）
% A5判。章内で定義・定理・例・注意等の番号を共通化。
% コンパイル： platex -interaction=nonstopmode math_note_template_v2.tex (2回)
%              dvipdfmx math_note_template_v2.dvi
\documentclass[a5paper,10pt,onecolumn]{jsarticle}
\usepackage[dvipdfmx,top=22mm,bottom=20mm,inner=19mm,outer=18mm,headheight=13pt,headsep=6mm]{geometry}
\usepackage{amsmath,amssymb,amsthm,mathtools}
\usepackage{enumitem,fancyhdr,titlesec}
\usepackage[dvipdfmx,bookmarks=true,colorlinks=false]{hyperref}
\usepackage{pxjahyper}

% ----- 差し替える文書情報 -----
\newcommand{\NoteTitle}{Fourier変換とPlancherelの定理}
\newcommand{\NoteSeries}{数学ノート集}

% ----- 文字・組版 -----
% 参考PDFの質素な明朝＋Computer Modernの印象を引き継ぎながら，
% A5判で窮屈にならないよう本文の行間を少し確保する．
\linespread{1.13}
\setlength{\parindent}{1zw}
\setlength{\parskip}{0pt}
\setlength{\abovedisplayskip}{6pt plus 2pt minus 2pt}
\setlength{\belowdisplayskip}{6pt plus 2pt minus 2pt}
\setlength{\abovedisplayshortskip}{3pt plus 2pt minus 1pt}
\setlength{\belowdisplayshortskip}{4pt plus 2pt minus 1pt}
\setlist[enumerate]{topsep=3pt,itemsep=1pt,parsep=0pt,leftmargin=2.4zw}
\titleformat{\section}{\normalfont\normalsize\bfseries}{\thesection}{0.8zw}{}
\titlespacing*{\section}{0pt}{2.0ex plus .4ex}{.9ex}
\titleformat{\subsection}{\normalfont\normalsize\bfseries}{\thesubsection}{0.7zw}{}
\titlespacing*{\subsection}{0pt}{1.5ex plus .3ex}{.6ex}
\raggedbottom

% ----- 柱とページ番号：ページ上部、罫線なし -----
\fancypagestyle{mathnote}{%
  \fancyhf{}%
  \fancyhead[L]{\footnotesize \NoteSeries\ （\NoteTitle）}%
  \fancyhead[R]{\footnotesize \thepage}%
  \renewcommand{\headrulewidth}{0pt}%
  \renewcommand{\footrulewidth}{0pt}%
}
\pagestyle{mathnote}

% ----- 節内通し番号の定義・定理群 -----
\newtheoremstyle{mathnotetheorem}
 {4.5pt}{4.5pt}{\normalfont}{}%
 {\bfseries}{}{0.65zw}%
 {\thmname{#1}\thmnumber{ #2}\thmnote{（#3）}}
\theoremstyle{mathnotetheorem}
\newtheorem{theorem}{定理}[section]
\newtheorem{definition}[theorem]{定義}
\newtheorem{lemma}[theorem]{補題}
\newtheorem{proposition}[theorem]{命題}
\newtheorem{corollary}[theorem]{系}
\newtheorem{example}[theorem]{例}
\newtheorem{remark}[theorem]{注意}
\renewcommand{\proofname}{証明}
\renewcommand{\qedsymbol}{$\square$}
\makeatletter
\renewenvironment{proof}[1][\proofname]{%
  \par\pushQED{\qed}\normalfont\topsep4\p@\@plus2\p@\relax
  \trivlist\item[\hskip\labelsep\bfseries #1]\ignorespaces
}{\popQED\endtrivlist\@endpefalse}
\makeatother

% ----- 使うノートで適宜拡張する記号 -----
\newcommand{\R}{\mathbb{R}}
\newcommand{\norm}[1]{\lVert #1\rVert}
\newcommand{\abs}[1]{\lvert #1\rvert}

\begin{document}
\section*{読了後にできること}
本稿を読むための前提は，Fourier変換，Plancherelの定理，Gaussian積分．
本稿の具体的な到達課題は，具体的な関数のFourier変換を積分から求める．


\par\noindent\textbf{記号と約束.}\quad
関数空間の添字 $c$ はコンパクトな台をもつことを表す．
\par

\section{$L^1$ 上のFourier変換}

本稿では $\mathbb R^d$ 上のLebesgue積分を用いる．
Fourier変換の正規化を
\[
 \widehat f(\xi)=\mathcal Ff(\xi)
 =(2\pi)^{-d/2}\int_{\mathbb R^d}
 f(x)e^{-ix\cdot\xi}\,dx
\]
と定める．

\begin{proposition}[基本的な有界性]\label{prop:l1-bounded}
$f\in L^1(\mathbb R^d)$ なら
$\widehat f$ は有界連続関数であり，
\[
 \norm{\widehat f}_\infty
 \le(2\pi)^{-d/2}\norm f_1.
\]
\end{proposition}
\begin{proof}
三角不等式から
\[
 |\widehat f(\xi)|
 \le(2\pi)^{-d/2}\int_{\mathbb R^d}|f(x)|\,dx
\]
を得る．
$\xi_n\to\xi$ なら
$f(x)e^{-ix\cdot\xi_n}\to f(x)e^{-ix\cdot\xi}$
が各 $x$ で成立し，
その絶対値は $|f(x)|$ で抑えられる．
Lebesgueの収束定理によって
$\widehat f(\xi_n)\to\widehat f(\xi)$ となり連続である．
\end{proof}

\begin{proposition}[変換公式]\label{prop:basic-transforms}
$f\in L^1(\mathbb R^d)$，
$a,b\in\mathbb R^d$，$\lambda\in\mathbb R\setminus\{0\}$
とする．このとき
\begin{align*}
 \mathcal F(f(\cdot-a))(\xi)&=e^{-ia\cdot\xi}\widehat f(\xi),\\
 \mathcal F(e^{ib\cdot x}f(x))(\xi)&=\widehat f(\xi-b),\\
 \mathcal F(f(\lambda\cdot))(\xi)
 &=|\lambda|^{-d}\widehat f(\xi/\lambda).
\end{align*}
\end{proposition}
\begin{proof}
第一式は変数変換 $y=x-a$ により
\[
 \int f(x-a)e^{-ix\cdot\xi}\,dx
 =e^{-ia\cdot\xi}\int f(y)e^{-iy\cdot\xi}\,dy
\]
となることから従う．
第二式は指数関数をまとめて
$e^{ib\cdot x}e^{-ix\cdot\xi}=e^{-ix\cdot(\xi-b)}$
とすればよい．
第三式は $y=\lambda x$ と置く変数変換で，
Jacobianの絶対値が $|\lambda|^{-d}$ であることによる．
\end{proof}

\section{Schwartz空間}

\begin{definition}
$\mathcal S(\mathbb R^d)$ を，無限回微分可能で，
任意の多重指数 $\alpha,\beta$ について
\[
 \sup_{x\in\mathbb R^d}
 |x^\alpha\partial^\beta f(x)|<\infty
\]
を満たす関数全体とする．
これをSchwartz空間という．
\end{definition}

\begin{lemma}\label{lem:schwartz-l1}
任意の $f\in\mathcal S(\mathbb R^d)$ と
多重指数 $\alpha,\beta$ に対して
$x^\alpha\partial^\beta f\in L^1(\mathbb R^d)$ である．
\end{lemma}
\begin{proof}
$N>d$ を取る．
Schwartz空間の定義から
\[
 (1+|x|)^N|x^\alpha\partial^\beta f(x)|
 \le C
\]
となる定数 $C$ が存在する．
従って
$|x^\alpha\partial^\beta f(x)|
 \le C(1+|x|)^{-N}$ である．
$N>d$ なので右辺は $\mathbb R^d$ で可積分であり，
結論を得る．
\end{proof}

\begin{theorem}[Schwartz空間のFourier不変性]\label{thm:schwartz}
$f\in\mathcal S(\mathbb R^d)$ なら
$\widehat f\in\mathcal S(\mathbb R^d)$ である．
さらに
\[
 \mathcal F(\partial_jf)(\xi)=i\xi_j\widehat f(\xi),
 \qquad
 \partial_{\xi_j}\widehat f(\xi)
 =-i\mathcal F(x_jf)(\xi).
\]
\end{theorem}
\begin{proof}
$f$ と $x_jf$ は可積分なので，
積分と $\xi_j$ に関する微分を交換して
\[
 \partial_{\xi_j}\widehat f(\xi)
 =(2\pi)^{-d/2}\int_{\mathbb R^d}
 (-ix_j)f(x)e^{-ix\cdot\xi}\,dx
 =-i\mathcal F(x_jf)(\xi)
\]
を得る．
一方，$f$ は急減少するので
$|x|\to\infty$ の境界項は零となる．
$x_j$ に関して部分積分すると
\[
 \mathcal F(\partial_jf)(\xi)
 =(2\pi)^{-d/2}\int
 \partial_jf(x)e^{-ix\cdot\xi}\,dx
 =i\xi_j\widehat f(\xi)
\]
となる．
これらを繰り返すと
$\xi^\alpha\partial_\xi^\beta\widehat f(\xi)$ は
Schwartz関数の有限個の多項式倍・微分の
Fourier変換の定数倍として表せる．
補題\ref{lem:schwartz-l1}と
命題\ref{prop:l1-bounded}によって
それらは一様有界である．
従って $\widehat f\in\mathcal S$ となる．
\end{proof}

\section{Gaussian関数のFourier変換}

\begin{lemma}[Gaussian積分]\label{lem:gaussian-integral}
\[
 \int_{\mathbb R}e^{-x^2/2}\,dx=\sqrt{2\pi}.
\]
\end{lemma}
\begin{proof}
$I=\int_{\mathbb R}e^{-x^2/2}\,dx$ と置く．
Tonelliの定理によって
\[
 I^2=\int_{\mathbb R^2}e^{-(x^2+y^2)/2}\,dx\,dy.
\]
極座標変換を適用すると
\[
 I^2=\int_0^{2\pi}\int_0^\infty
 e^{-r^2/2}r\,dr\,d\theta
 =2\pi.
\]
$I>0$ なので $I=\sqrt{2\pi}$ である．
\end{proof}

\begin{theorem}[Gaussianの自己Fourier性]\label{thm:gaussian}
$G(x)=e^{-|x|^2/2}$ と置く．このとき
\[
 \widehat G(\xi)=e^{-|\xi|^2/2}.
\]
\end{theorem}
\begin{proof}
まず $d=1$ とする．$G'(x)=-xG(x)$ なので
定理\ref{thm:schwartz}から
\[
 \widehat G'(\xi)
 =-i\mathcal F(xG)(\xi)
 =i\mathcal F(G')(\xi)
 =i(i\xi)\widehat G(\xi)
 =-\xi\widehat G(\xi).
\]
初期値は補題\ref{lem:gaussian-integral}によって
$\widehat G(0)=1$ である．
常微分方程式 $h'=-\xi h$ の一意な解は
$h(\xi)=e^{-\xi^2/2}$ なので結論を得る．
一般の $d$ については
$G(x)=\prod_{j=1}^de^{-x_j^2/2}$ であり，
Fubiniの定理で各座標の積分に分解すれば
$\widehat G(\xi)=\prod_j e^{-\xi_j^2/2}
=e^{-|\xi|^2/2}$ となる．
\end{proof}

\section{Fourier反転公式}

\begin{lemma}[Gaussian近似恒等作用素]\label{lem:gaussian-approx}
$f\in\mathcal S(\mathbb R^d)$ に対し
\[
 \rho_\varepsilon(x)
 =(2\pi\varepsilon)^{-d/2}e^{-|x|^2/(2\varepsilon)}
 \qquad(\varepsilon>0)
\]
と置く．このとき
\[
 (\rho_\varepsilon*f)(x)\to f(x)
 \qquad(\varepsilon\downarrow0)
\]
が各 $x\in\mathbb R^d$ で成り立つ．
\end{lemma}
\begin{proof}
Gaussian積分と変数変換により
$\int\rho_\varepsilon=1$ である．
\[
 (\rho_\varepsilon*f)(x)-f(x)
 =\int_{\mathbb R^d}
 \rho_\varepsilon(y)(f(x-y)-f(x))\,dy.
\]
任意の $\delta>0$ に対して積分領域を
$|y|<\delta$ と $|y|\ge\delta$ に分ける．
$f$ は連続なので，前者では $\delta$ を十分小さく取れば
$|f(x-y)-f(x)|<\eta$ となる．
後者では
$|f(x-y)-f(x)|\le2\norm f_\infty$ であり，
Gaussian変数変換 $y=\sqrt\varepsilon z$ によって
$\int_{|y|\ge\delta}\rho_\varepsilon(y)dy
 \to0$ となる．
従って極限は零である．
\end{proof}

\begin{theorem}[Fourier反転公式]\label{thm:inversion}
$f\in\mathcal S(\mathbb R^d)$ に対して
\[
 f(x)=(2\pi)^{-d/2}
 \int_{\mathbb R^d}\widehat f(\xi)e^{ix\cdot\xi}\,d\xi
\]
が全ての $x$ で成り立つ．
特に
\[
 \mathcal F^2f(x)=f(-x).
\]
\end{theorem}
\begin{proof}
$\varepsilon>0$ とし，
\[
 I_\varepsilon(x)
 =(2\pi)^{-d/2}\int_{\mathbb R^d}
 \widehat f(\xi)e^{ix\cdot\xi}
 e^{-\varepsilon|\xi|^2/2}\,d\xi
\]
と置く．$\widehat f$ はSchwartz関数なので
積分は絶対収束する．
$\widehat f$ の定義を代入すると
\begin{align*}
 I_\varepsilon(x)
 &=(2\pi)^{-d}\int_{\mathbb R^d}\int_{\mathbb R^d}
 f(y)e^{i(x-y)\cdot\xi}
 e^{-\varepsilon|\xi|^2/2}\,dy\,d\xi.
\end{align*}
$f\in L^1$ でGaussian核も $L^1$ なので
Fubiniの定理が適用できる．
定理\ref{thm:gaussian}と拡大縮小の公式から
\[
 (2\pi)^{-d}\int_{\mathbb R^d}
 e^{i(x-y)\cdot\xi}e^{-\varepsilon|\xi|^2/2}\,d\xi
 =\rho_\varepsilon(x-y)
\]
である．従って $I_\varepsilon(x)=(\rho_\varepsilon*f)(x)$
となり，補題\ref{lem:gaussian-approx}から
$I_\varepsilon(x)\to f(x)$ を得る．
一方 $\widehat f\in L^1$ なので，
Lebesgueの収束定理によって
\[
 I_\varepsilon(x)\to
 (2\pi)^{-d/2}\int\widehat f(\xi)e^{ix\cdot\xi}\,d\xi.
\]
両極限を比較して反転公式を得る．
最後の式は反転公式の右辺が
$\mathcal F(\widehat f)(-x)$ に等しいことによる．
\end{proof}

\section{Parsevalの等式とPlancherelの定理}

\begin{proposition}[Parsevalの恒等式]\label{prop:parseval}
$f,g\in\mathcal S(\mathbb R^d)$ に対し
\[
 \int_{\mathbb R^d}f(x)\overline{g(x)}\,dx
 =\int_{\mathbb R^d}
   \widehat f(\xi)\overline{\widehat g(\xi)}\,d\xi
\]
が成り立つ．特に $\norm{\widehat f}_2=\norm f_2$ である．
\end{proposition}
\begin{proof}
反転公式から
\[
 \overline{g(x)}
 =(2\pi)^{-d/2}
 \int_{\mathbb R^d}
 \overline{\widehat g(\xi)}e^{-ix\cdot\xi}\,d\xi
\]
である．これを左辺の積分に代入する．
$f,\widehat g$ はともにSchwartz関数で $L^1$ に属するから
Fubiniの定理が適用でき，
\begin{align*}
 \int f(x)\overline{g(x)}\,dx
 &=(2\pi)^{-d/2}\int\int
 f(x)e^{-ix\cdot\xi}
 \overline{\widehat g(\xi)}\,dx\,d\xi\\
 &=\int\widehat f(\xi)\overline{\widehat g(\xi)}\,d\xi.
\end{align*}
$f=g$ と置いてノルムの一致を得る．
\end{proof}

\begin{lemma}[Schwartz関数の稠密性]\label{lem:schwartz-dense}
$\mathcal S(\mathbb R^d)$ は $L^2(\mathbb R^d)$ で稠密である．
\end{lemma}
\begin{proof}
Lebesgueの収束定理による切断
$f1_{\{|x|\le R,\ |f|\le R\}}$ により，
任意の $L^2$ 関数は有界なコンパクト台関数で近似できる．
Lebesgue測度の正則性を用いて，
そのような関数はコンパクト台連続関数で
$L^2$ 近似できる．
さらにコンパクト台連続関数を
滑らかなコンパクト台の近似恒等作用素と畳み込むと，
$C_c^\infty$ 関数で一様近似できる．
$C_c^\infty\subset\mathcal S$ なので
$\mathcal S$ は $L^2$ で稠密である．
\end{proof}

\begin{theorem}[Plancherelの定理]\label{thm:plancherel}
Fourier変換
$\mathcal F:\mathcal S(\mathbb R^d)\to\mathcal S(\mathbb R^d)$
は，一意な線形等距離全単射
\[
 \mathcal F:L^2(\mathbb R^d)\to L^2(\mathbb R^d)
\]
に拡張される．
この拡張はユニタリ作用素であり，
\[
 \norm{\mathcal Ff}_2=\norm f_2,\qquad
 \mathcal F^{-1}f(x)=\mathcal Ff(-x)
\]
が $L^2$ の意味で成り立つ．
\end{theorem}
\begin{proof}
$f\in L^2$ に対し，
補題\ref{lem:schwartz-dense}によって
$\norm{f_n-f}_2\to0$ となる
$f_n\in\mathcal S$ を取る．
Parsevalの恒等式により
\[
 \norm{\widehat f_n-\widehat f_m}_2
 =\norm{f_n-f_m}_2
\]
なので $(\widehat f_n)$ は $L^2$ のCauchy列である．
$L^2$ は完備だから極限が存在する．
この極限を $\mathcal Ff$ と定める．
別の近似列を取っても差の $L^2$ ノルムが零へ向かうので
定義は近似列によらない．
また線形性とノルムの等式は極限を取ることで保たれる．
従って $\mathcal F$ は等距離線形写像であり，
その像は閉である．
一方，反転公式によって
$\mathcal F(\mathcal S)=\mathcal S$ であり，
$\mathcal S$ は $L^2$ に稠密なので，
$\mathcal F$ の像は稠密でもある．
閉かつ稠密な像は $L^2$ 全体であるから
$\mathcal F$ は全射となる．
Parsevalの恒等式も極限へ移せるので内積を保存し，
従ってユニタリである．
最後に
$\mathcal F^2f(x)=f(-x)$ はSchwartz関数で成立し，
$L^2$ での近似と鏡映の等距離性によって
一般の $L^2$ 関数へ拡張できる．
\end{proof}

\section{畳み込みとFourier変換}

\begin{definition}
$f,g\in L^1(\mathbb R^d)$ に対して
\[
 (f*g)(x)=\int_{\mathbb R^d}f(y)g(x-y)\,dy
\]
を畳み込みという．
\end{definition}

\begin{proposition}[畳み込み定理]\label{prop:convolution}
$f,g\in L^1(\mathbb R^d)$ なら
$f*g\in L^1$ であり，
\[
 \widehat{f*g}(\xi)
 =(2\pi)^{d/2}\widehat f(\xi)\widehat g(\xi)
\]
が成り立つ．
\end{proposition}
\begin{proof}
Tonelliの定理によって
\[
 \int|f*g(x)|\,dx
 \le\int\int|f(y)||g(x-y)|\,dy\,dx
 =\norm f_1\norm g_1
\]
である．従って $f*g\in L^1$ となる．
Fourier変換の定義に代入して
Fubiniの定理を適用すると
\begin{align*}
 \widehat{f*g}(\xi)
 &=(2\pi)^{-d/2}\int\int
 f(y)g(x-y)e^{-ix\cdot\xi}\,dy\,dx\\
 &=(2\pi)^{-d/2}
 \left(\int f(y)e^{-iy\cdot\xi}dy\right)
 \left(\int g(z)e^{-iz\cdot\xi}dz\right)\\
 &=(2\pi)^{d/2}\widehat f(\xi)\widehat g(\xi).
\end{align*}
\end{proof}

\section{熱方程式への応用}

\begin{theorem}[熱方程式のFourier解]\label{thm:heat}
$f\in\mathcal S(\mathbb R^d)$ とし，
\[
 u(t,x)=(4\pi t)^{-d/2}
 \int_{\mathbb R^d}
 e^{-|x-y|^2/(4t)}f(y)\,dy,\qquad t>0
\]
と置く．このとき $u$ は $t>0$ で滑らかであり，
\[
 \partial_tu=\Delta u,\qquad
 u(t,\cdot)\to f
\]
が $t\downarrow0$ で $L^2$ の意味で成り立つ．
また
\[
 \widehat{u(t,\cdot)}(\xi)
 =e^{-t|\xi|^2}\widehat f(\xi)
\]
が成り立つ．
\end{theorem}
\begin{proof}
$k_t(x)=(4\pi t)^{-d/2}e^{-|x|^2/(4t)}$ と置くと
$u(t,\cdot)=k_t*f$ である．
Gaussian変換の公式から
\[
 \widehat{k_t}(\xi)
 =(2\pi)^{-d/2}e^{-t|\xi|^2}.
\]
従って畳み込み定理によって
$\widehat u=e^{-t|\xi|^2}\widehat f$ を得る．
両辺を $t$ で微分すると
$\partial_t\widehat u=-|\xi|^2\widehat u$ である．
一方，Schwartz関数上の微分公式により
$\widehat{\Delta u}=-|\xi|^2\widehat u$ なので，
Fourier反転公式から $\partial_tu=\Delta u$ となる．
また $t\downarrow0$ では
$e^{-t|\xi|^2}\widehat f(\xi)\to\widehat f(\xi)$ であり，
絶対値が $|\widehat f(\xi)|$ 以下なので，
$L^2$ のLebesgue収束定理から
\[
 \norm{(e^{-t|\xi|^2}-1)\widehat f}_2\to0.
\]
Plancherelの定理によって
$\norm{u(t,\cdot)-f}_2\to0$ となる．
\end{proof}

\section{不確定性不等式}

\begin{theorem}[Heisenbergの不確定性不等式]\label{thm:heisenberg}
任意の $f\in\mathcal S(\mathbb R^d)$ に対して
\[
 \left(\int_{\mathbb R^d}|x|^2|f(x)|^2\,dx\right)
 \left(\int_{\mathbb R^d}|\xi|^2|\widehat f(\xi)|^2\,d\xi\right)
 \ge\frac{d^2}{4}\norm f_2^4.
\]
\end{theorem}
\begin{proof}
$f$ は急減少するので
$\int_{\mathbb R^d}\operatorname{div}(x|f(x)|^2)\,dx=0$
である．
発散を展開すると
\[
 0=d\norm f_2^2
 +2\operatorname{Re}\int_{\mathbb R^d}
 (x\cdot\nabla f(x))\overline{f(x)}\,dx.
\]
従ってCauchy--Schwarzの不等式から
\[
 d\norm f_2^2
 \le2\left(\int|x|^2|f|^2\right)^{1/2}
       \left(\int|\nabla f|^2\right)^{1/2}.
\]
Fourier変換の微分公式とPlancherelの定理によって
\[
 \int|\nabla f(x)|^2dx
 =\int|\xi|^2|\widehat f(\xi)|^2d\xi
\]
となる．これを前の不等式に代入して両辺を二乗すれば
結論を得る．
\end{proof}

\begin{remark}
Fourier変換は空間の平行移動や微分を，
周波数空間の乗法へ変換する．
Plancherelの定理はこの対応を
$L^2$ 空間上のユニタリ同型として厳密に保証する．
\end{remark}

\section{到達課題の完全解答}
\noindent\textbf{Gaussian関数のFourier変換．}
\begin{proposition}
Fourier変換を
\[
\widehat f(\xi)=\int_{\mathbb R}f(x)e^{-2\pi ix\xi}\,dx
\]
と定めると，
$f(x)=e^{-\pi x^2}$ に対して
$\widehat f(\xi)=e^{-\pi\xi^2}$ である．
\end{proposition}
\begin{proof}
$F(\xi)=\int_{\mathbb R}
e^{-\pi x^2}e^{-2\pi ix\xi}\,dx$ と置く．
積分と微分を交換し，部分積分すると
\[
\begin{aligned}
F'(\xi)
&=-2\pi i\int_{\mathbb R}
xe^{-\pi x^2}e^{-2\pi ix\xi}\,dx\\
&=i\int_{\mathbb R}
\frac d{dx}(e^{-\pi x^2})
e^{-2\pi ix\xi}\,dx
=-2\pi\xi F(\xi).
\end{aligned}
\]
従って $F(\xi)=Ce^{-\pi\xi^2}$．
Gaussian積分から $F(0)=1$ なので $C=1$．
\end{proof}
\noindent\textbf{到達確認．}
$f(x)=e^{-\pi ax^2}$（$a>0$）のFourier変換を求めよ．
\par\noindent\textbf{解答．}
変数変換 $y=\sqrt a\,x$ より
$\widehat f(\xi)=a^{-1/2}e^{-\pi\xi^2/a}$ となる．
\end{document}